% TOP_PROBLEM: {"id": 20002556, "title": "Exact reversible shooting criterion for halo-orbit bifurcation", "category_id": 16, "category": {"id": 16, "name": "physics", "display_name": "Mathematical Physics", "description": "Problems at the intersection of mathematics and physics.", "slug": "physics", "order_index": 16, "created_at": "2026-07-31T15:26:25.671Z"}, "status": "partially_solved", "statusQualification": "The primary source includes a tentative bifurcation remark and states that the author knew numerical studies. That historical remark is not independently checked or promoted to a precise theorem. Parameters and family definition must be specified for a proof. The catalogue title/status was paired with a machine-generated research attempt; that attempt is not peer-reviewed evidence of settlement. The statement covers the original source questions rather than only the title’s subcase."}
% FIRST_POSTED: 2026-10-11
% ENTRY_KIND: statement_only
% SUBMISSION: {"kind": "statement_only", "category": "top_problems", "problem_id": "20002556", "proposal_type": "established", "source": "American Institute of Mathematics, Variational Methods in Celestial Mechanics; notes by Kuo-Chang Chen, version 22 June 2003, source item 27. https://aimath.org/WWN/varcelest/varcelest.pdf", "status_evidence": "Inherited source-catalogue status partially_solved preserved. Background has machine-generated research claims explicitly marked not peer reviewed; these are not independent status evidence. Primary source accessed 11 October 2026; compared against the supplied catalogue excerpt. Historical wording and partial-result remarks were checked for scope; no comprehensive current-literature or proof-verification audit was performed.", "status_checked_on": "2026-10-11", "duplicate_check": "Full 15,773-record source registry checked for identical normalized statements; available site registry checked for identical normalized title/statement and canonical-ID collision. Canonical ID is absent from the site registry. No mathematical-equivalence literature audit performed.", "attribution": "Rick Moeckel. Statement preparation: OpenAI GPT-6 Codex assistant, generation 11 October 2026; source-matching and notation editing by the assistant, no independent human mathematical verification."}
% SOURCE_STATEMENT: "Problem 27. Existence proof for the \"halo orbits\" of the three-dimensional restricted three-body problem. (Rick Moeckel)"
% Preparation and final self-audit: CONTRIBUTING.md prompts C and D; source/audit artifacts retained in the local physics_statements_2026_10_11 batch.
% Imported from: physics_statements.tex; UnsolvedMath ID 20002556
% !TeX program = lualatex
\documentclass[11pt]{article}
\usepackage[margin=25mm]{geometry}
\usepackage{fontspec}
\setmainfont{Latin Modern Roman}
\usepackage{amsmath,amssymb,mathtools}
\usepackage{unicode-math}
\setmathfont{Latin Modern Math}
\usepackage{xurl,hyperref}
\hypersetup{colorlinks=true,urlcolor=blue}
\setcounter{secnumdepth}{0}
\setlength{\parindent}{0pt}
\setlength{\parskip}{5pt}
\setlength{\emergencystretch}{3em}
\newcommand{\sref}[2]{\hyperlink{source:#1:#2}{[S#2]}}
\newcommand{\cataloguescope}[1]{\paragraph{Recorded target.}#1}
\newcommand{\R}{\mathbb R}
\newcommand{\C}{\mathbb C}
\newcommand{\N}{\mathbb N}
\newcommand{\Z}{\mathbb Z}
\newcommand{\Q}{\mathbb Q}
\newcommand{\F}{\mathbb F}
\newcommand{\eps}{\varepsilon}
\begin{document}
% UnsolvedMath ID 20002556; source catalogue code AIM-PHYSICS-0064
\setcounter{section}{20002555}
\section{Exact reversible shooting criterion for halo-orbit bifurcation}\label{20002556}
{\small UnsolvedMath ID 20002556 · Mathematical Physics}
\subsection{Definitions and mathematical statement}

\paragraph{Shared notation.}$\mathbb N=\{1,2,\ldots\}$; the standard number fields have their usual meanings. All additional conventions are stated locally.


\paragraph{Definitions and notation.}In the circular restricted three-body problem, two primaries follow prescribed circular motion and a test particle has negligible effect on them. In rotating coordinates the Jacobi constant combines kinetic energy and the effective gravitational/centrifugal potential. The Hill region is the set where the zero-velocity inequality permits motion. Its necks occur near collinear equilibria. In the spatial restricted problem, halo orbits are the named three-dimensional periodic families near the collinear libration-point region, with characteristic symmetries and out-of-plane motion. Bifurcation from a planar periodic family and existence of a numerically observed family are different assertions.

\paragraph{Catalogue wording retained for traceability.}
Problem 27. Existence proof for the "halo orbits" of the three-dimensional restricted three-body problem. (Rick Moeckel)
\paragraph{Scope and source qualification.}The primary source includes a tentative bifurcation remark and states that the author knew numerical studies. That historical remark is not independently checked or promoted to a precise theorem. Parameters and family definition must be specified for a proof. The catalogue title/status was paired with a machine-generated research attempt; that attempt is not peer-reviewed evidence of settlement. The statement covers the original source questions rather than only the title’s subcase.
\subsection{Short English statement}
Give a rigorous existence proof for the spatial restricted three-body halo-orbit family.
\subsection{Sources}
\begin{itemize}
\item[\hypertarget{source:20002556:1}{S1}] ULAM.AI and the UnsolvedMath Contributors, UnsolvedMath, record 20002556 (AIM-PHYSICS-0064). Catalogue curation under CC BY 4.0; definitions and English exposition added locally.
\url{https://huggingface.co/datasets/ulamai/UnsolvedMath}.
\item[\hypertarget{source:20002556:2}{S2}] American Institute of Mathematics, Variational Methods in Celestial Mechanics; notes by Kuo-Chang Chen, version 22 June 2003, source item 27.
\url{https://aimath.org/WWN/varcelest/varcelest.pdf}.
\end{itemize}
\label{end:20002556}
\end{document}
